\chapter{Physics-Informed Neural Networks}
\label{ch:pinns}

\epigraph{\itshape A PINN is a neural network whose loss function is a
PDE.}{Raissi, Perdikaris, Karniadakis (2019)}

\section{Formulation}
Consider a PDE on $\Omega\times[0,T]$:
\begin{align}
\Nop[u](\bm{x},t)&=0 \quad \text{in } \Omega\times(0,T],\\
\Bop[u](\bm{x},t)&=g(\bm{x},t) \quad \text{on } \partial\Omega\times[0,T],\\
u(\bm{x},0)&=u_0(\bm{x}) \quad \text{in } \Omega.
\end{align}
A PINN parameterises $u\approx u_\theta:\R^{d+1}\to\R^k$ as a multilayer
perceptron and minimises a composite loss~\cite{raissi2019pinn}
\begin{equation}
\Lop(\theta) = w_r\Lop_r + w_b\Lop_b + w_0\Lop_0 + w_d\Lop_d,
\end{equation}
where
\begin{align}
\Lop_r &=\frac{1}{N_r}\sum_i \big|\Nop[u_\theta](\bm{x}_i,t_i)\big|^2,\\
\Lop_b &=\frac{1}{N_b}\sum_i \big|\Bop[u_\theta]-g\big|^2,\\
\Lop_0 &=\frac{1}{N_0}\sum_i \big|u_\theta(\bm{x}_i,0)-u_0\big|^2,\\
\Lop_d &=\frac{1}{N_d}\sum_i \big|u_\theta(\bm{x}_i,t_i)-u_i^{\text{obs}}\big|^2.
\end{align}
Derivatives appearing in $\Nop$ are obtained by automatic
differentiation through the network.

\section{Why it works (when it works)}
The PINN loss imposes the strong-form PDE on a finite set of
\emph{collocation points}. With sufficiently expressive $u_\theta$ and
adequate optimisation, $u_\theta$ approaches a strong solution. The
formulation is mesh-free, accommodates irregular geometries, and
unifies forward and inverse problems --- inverse coefficients become
trainable parameters.

\section{Variants}
\paragraph{XPINN / cPINN.}
Domain decomposition: separate networks per subdomain stitched together
by interface losses, enabling parallelism and capture of localised
features~\cite{jagtap2020xpinn}.

\paragraph{hp-VPINN.}
Use the variational (weak) form with test functions, integrating by
parts to reduce derivative orders --- helpful when solutions have low
regularity.

\paragraph{Causal PINN.}
Re-weight the time-domain residual to enforce that the loss at later
times is small only \emph{after} earlier times converge, avoiding
non-causal solutions to time-evolution PDEs~\cite{wang2022causal}.

\paragraph{NTK-based PINN.}
The neural tangent kernel reveals spectral bias in PINN training:
high-frequency residual modes converge slowly. Re-weighting based on
the NTK eigenvalues fixes the imbalance~\cite{wang2022ntkpinn}.

\paragraph{Self-adaptive / SA-PINN.}
Per-point trainable weights $\lambda_i$ adversarially upweight hard
collocation points~\cite{mcclenny2023sapinn}.

\paragraph{PINN-LS, gPINN, PIELM.}
Least-squares formulations, gradient-enhanced PINNs, and physics-informed
extreme learning machines trade off flexibility for speed.

\section{Failure modes}
\label{sec:pinn-failure}

PINNs are notoriously brittle. Documented failure modes include:
\begin{enumerate}
\item \textbf{Spectral bias} towards low frequencies; sharp gradients
and shock fronts under-resolved.
\item \textbf{Multi-objective imbalance}: residual, boundary, and data
losses conflict; gradient pathologies emerge.
\item \textbf{Stiffness in time}: optimiser converges to spurious steady
states matching the residual at $t=0$ but not propagating.
\item \textbf{Non-convex landscapes}: local minima where residual is
small but solution is wrong.
\item \textbf{Floating-point precision}: a 2025 result shows many
``failure modes'' disappear when training is performed in
\textsf{float64} rather than \textsf{float32}, because the L-BFGS
convergence test is otherwise triggered prematurely~\cite{rethinking2025pinn}.
\end{enumerate}

\begin{pitfall}{The ``trivial solution'' trap}
For homogeneous PDEs with zero forcing and inadequate boundary loss
weighting, a PINN can converge to $u_\theta\equiv 0$. Always sanity-check
the solution magnitude and inspect residual histograms by region.
\end{pitfall}

\subsection{Mitigations}
\begin{recipe}{Hardening a PINN}
\begin{enumerate}
\item Use \textsf{float64}; verify L-BFGS convergence flags.
\item Adaptive collocation: resample where residual is large (RAR,
RAD).
\item Loss balancing: NTK weights, gradient-norm balancing,
SoftAdapt, ReLoBRaLo.
\item Hard-constrain BCs/ICs by output transforms
$u_\theta(\bm{x},t) = u_0(\bm{x}) + t\,\bm{N}_\theta$.
\item Causal time-marching: train on expanding time windows.
\item Fourier features / SIREN activations to combat spectral bias.
\item Exact gradient computation; avoid \texttt{create\_graph=True}
pitfalls in PyTorch.
\end{enumerate}
\end{recipe}

\section{Inverse problems with PINNs}
PINNs shine in data-assimilation regimes where partial measurements
constrain unknown coefficients $\mu$ in $\Nop_\mu$. Joint optimisation
of $(\theta,\mu)$ gives forward + inverse for free. This is why PINNs
are popular in subsurface flow, biomechanics, and astrophysical fluid
inference.

\section{Theoretical guarantees}
Universal approximation extends to PINNs under mild regularity, but
\emph{trainability} is the binding constraint. Recent work bounds the
generalisation error in terms of the residual loss, neural tangent
kernel conditioning, and quadrature error of the collocation
sampling~\cite{wang2022ntkpinn,deryck2024pinns}.

\section{When to choose a PINN over alternatives}
\begin{insight}{PINNs vs. neural operators vs. classical solvers}
\begin{itemize}
\item Single problem instance, smooth solution, sparse data, want
mesh-free + inverse parameters: \emph{PINN}.
\item Many problem instances, repeated queries: \emph{neural operator}
(\cref{ch:operators}).
\item One-shot accuracy, complex geometry, no parameter unknowns:
classical FE/FV solver.
\item Stiff multiscale + partial data: \emph{universal differential
equations} (\cref{ch:hybrid-ude}).
\end{itemize}
\end{insight}

\section{Code skeleton (JAX)}
\begin{lstlisting}[language=Python,caption={Minimal PINN for 1D Burgers},label=lst:pinn]
import jax, jax.numpy as jnp, optax
from flax import linen as nn

class MLP(nn.Module):
    feats: tuple = (64,)*5
    @nn.compact
    def __call__(self, x):
        for f in self.feats:
            x = nn.tanh(nn.Dense(f)(x))
        return nn.Dense(1)(x)

def u_fn(params, x, t):
    return model.apply(params, jnp.stack([x, t]))[0]

def residual(params, x, t, nu=0.01/jnp.pi):
    u   = u_fn(params, x, t)
    u_t = jax.grad(u_fn, argnums=2)(params, x, t)
    u_x = jax.grad(u_fn, argnums=1)(params, x, t)
    u_xx= jax.grad(lambda p,x,t: jax.grad(u_fn,argnums=1)(p,x,t),argnums=1)(params,x,t)
    return u_t + u*u_x - nu*u_xx
\end{lstlisting}

\section{Summary}
PINNs remain a swiss-army knife when the equation is known, geometry is
awkward, and data is scarce. Their reputation for fragility is partly
fair and partly due to under-engineered training. With float64,
adaptive sampling, balanced losses, and causal training, vanilla PINNs
solve a remarkable range of problems --- a fact that the literature has
recently re-affirmed~\cite{rethinking2025pinn}.
